% =====================================================================
%  Section 01 — Layla's Story (verbatim user paraphrase)
% =====================================================================
\makesectioncover{Layla's Cairo}

\section{Layla's Story}
\label{sec:layla}

\subsection{6:15 AM, Imbaba}

It is early in the morning (6:15) and Layla, who lives in Imbaba, is
already up and her neighbourhood is already loud. Imbaba, which is on
the west of the Nile, is considered the densest neighbourhood in Egypt
and some argue it has one of the highest densities in the entire world
(\stat{63k people per square kilometre}), which is much larger than
Manhattan for example (27k) and some other places in Cairo such as New
Cairo (3k).

It is important to note that no metro line existed in Imbaba until
2024, and no metro station exists in an acceptable walking distance
from Layla's house.

\begin{callout}
\textbf{\color{plotblue}Note.}\quad Layla is a fictional character that
represents the \stat{1.4 million Cairenes} that don't have formal
transport access. All the data used in this story is from notebooks,
datasets, and websites scraped: e.g. \texttt{citypopulation.de} for
the density of Imbaba, the \stat{83-minute trip time} is from
\texttt{explorecity.life}.
\end{callout}

\subsection{The office is in Maadi}

Layla's office is in Maadi (31.2581, 29.9601) which is 11 km away (on
paper) from where she lives (31.2083, 30.0866). Unfortunately, Layla
can't fly and has to use two transport systems that are not even
designed to be connected:

\begin{itemize}
  \item \kw{System A} (formal): the system built by the state, which
        is Cairo Metro, the new LRT, and the Ring-Road BRT.
  \item \kw{System B} (informal): the informal network including
        \stat{1{,}302} microbus stops, \stat{280} terminals (mawkf),
        and the knowledge every Cairene has in their head.
\end{itemize}

\subsection{The route is 83 minutes one way}

Layla has to take three vehicles, two transfers, and no one can plan
the trip for her end-to-end.

\begin{table}[H]
\centering\small
\renewcommand{\arraystretch}{1.25}
\begin{tabular}{@{} >{\color{plotorange}\bfseries}l
                    >{\color{bodytext}}l
                    >{\color{plotyellow}}r @{}}
\toprule
\color{plotblue}Leg & \color{plotblue}Mode & \color{plotblue}Minutes \\
\midrule
Walk to mawfaq         & on foot                  & 8  \\
Wait at mawfaq         & informal terminal queue  & 8.5 \\
Microbus 1             & informal microbus        & 25 \\
Walk between systems   & on foot                  & 12 \\
Metro Line 2 to Maadi  & formal metro             & 30 \\
\midrule
\textbf{Total one-way} &                          & \stat{83} \\
\bottomrule
\end{tabular}
\caption{Table 1. Layla spends around 166 minutes every day, which
results in Layla losing \stat{720 hours} every year just in
transportation.}
\end{table}

\subsection{At 6:15 it is not her choice, the mawfaq decides for her}

Five microbuses are loading customers at the same time, five
destinations written on cardboard, and drivers shouting route names.
Clearly there is no schedule, no app, and no driver cares about her
specific trip.

She then decides to choose a microbus. We can confirm that the mawfaq
she is at is not linked to any post-2014 metro stations. In short, the
investment of the last decade did not help Layla this morning.

\subsection{What is the problem?}

Cairo has two parallel transport systems and since 2014 more than
\stat{\$10 billion} were spent expanding System A (formal), which
includes Metro Line 3 (Ataba to Adly Mansour), the Cairo LRT to the
new administrative capital, and the Ring-Road BRT. However, the daily
experience of over \stat{1 million Cairenes} (Layla) has not changed
and did not become better.

\begin{callout}
\textbf{\color{plotblue}The big question: WHY?}\quad That is due to
the fact that System A was built where people \emph{will} live based
on state plans, however System B (informal) is carrying people where
they already live. In geography we will find both systems
overlapping, but in functionality? Not as much. Neither the formal
nor the informal has a timetable that references the other (no
formal transportation guide tells me what informal transportation I
should use, and vice versa), making it up to the poor Cairene who
has to plan his every-day route.
\end{callout}

Our report studies this disconnection with data. In Phase 1 we cleaned
\stat{7 GeoJSON datasets} which described the current reality of
transportation (\stat{1{,}302} stops, \stat{1{,}784} routes, \stat{280}
terminals, \stat{9{,}258} road segments, \stat{1{,}525} population
hexagons). We also answered \stat{12 research questions} about these
seven datasets and most importantly we found four structural gaps:
\stat{115 ghost terminals}, \stat{19 empty-return terminals}, vehicle
mismatch on \stat{75\% of long routes}, and also \stat{70 underserved
population hexagons}.

In Phase 2 we scraped \stat{eight new sources} which helped us
identify what the state has built since 2014 (\stat{89} metro
stations, \stat{20} LRT stations, \stat{21} BRT stations, \stat{68}
districts population, vehicle ownership, GTFS, OSM). We integrated
these scraped data with Phase 1 data by making a \stat{4-stage
pipeline} that goes as follows:

\begin{center}
KNN spatial $\to$ OSM verification $\to$ RapidFuzz fuzzy text
$\to$ SBERT multilingual matching
\end{center}

and then we asked \stat{13 questions} followed by testing
\stat{3 hypotheses} on whether the new infrastructure helped close
the gaps output in Phase 1 or not.

\begin{callout}
\textbf{\color{plotblue}What we found?}\quad We made a
\stat{3$\times$4} matrix that describes the percentage of how much we
closed the gaps after spending \stat{10 billion}, and the results were
not the best: \stat{no cell resulted in a fixing of above 25\%} and
most are below \stat{16\%}. In short, the new infrastructure did not
help Layla. The data showed us why and the report will walk us
through every step of how we came to that conclusion.
\end{callout}

\subsection{How this report is organised}

\begin{enumerate}
  \item \textbf{Phase 1:} what Cairo's existing transport system looks
        like. The data sources, data exploration, cleaning, KNN
        imputation, spatial merges, 12 research exploratory questions
        with their appropriate visualisation, and most importantly
        output 4 structural gaps.
  \item \textbf{Phase 2:} what did our state build after 2014. We
        scraped 8 sources and made a four-stage integration pipeline,
        and we also discussed 13 research questions with their
        appropriate visualisation.
  \item \textbf{Machine learning (AI) techniques:} we used 3 methods
        to assist us in this project (SBERT, K-Means, Moran's I); we
        chose them due to the fact that classical methods can't solve
        the same problem.
  \item \textbf{Hypothesis tests:} an explanation of all statistical
        tests used (Spearman, Kruskal, Mann-Whitney, Cliff's, epsilon
        squared, Moran's I, K-Means), then a clear walkthrough of
        \kw{H1} coverage-need mismatch, \kw{H2} which is LRT
        catchment deficit, and \kw{H3} which is BRT corridor match.
  \item \textbf{Synthesis:} a very important section on what all the
        evidence integrated together says, including the gap-closure
        matrix which is the most important result of this project,
        and the resolution set (how we reduced the time Layla
        consumes in transportation every day by half, from
        \stat{83 minutes} to \stat{40 minutes}), delivering a happy
        ending to Layla's sad story.
  \item \textbf{Appendices:} simple file manifest and glossary.
\end{enumerate}
